%% ======================================================================
\section{Mathematical toolbox}
\label{chap:tools}
%% ======================================================================

\begin{objectives}
\guiding{How can an error measured on a sample be turned into a guarantee on the
true error, first for one classifier, then for many?}
We fix the formal setting (Gibbs classifier and majority vote), and we gather
the probabilistic tools that are used over and over in PAC-Bayesian proofs:
concentration inequalities for a single voter, stated with the binary
Kullback--Leibler divergence, the Kullback--Leibler divergence between
distributions over voters, and the change of measure inequality, which is the
key step of every PAC-Bayesian bound. The section ends with the Occam bound,
which already contains the main idea of PAC-Bayes.
\end{objectives}

The section climbs three steps. We first learn how to certify \emph{one}
classifier fixed in advance (concentration inequalities). We then introduce the
tool that lets us move from a distribution fixed in advance to a distribution
chosen with the data (the change of measure). Finally, combining the two in the
simplest case gives the Occam bound, which certifies a classifier \emph{chosen}
among many. The next section will take the last step, from one chosen
classifier to a chosen distribution over classifiers.

\subsection{Setting}

Before stating any bound, we must fix the objects that the bounds talk about:
the data, the voters, their risks, and the two ways of combining voters.
We consider binary classification. Examples $(x,y)$ are drawn independently from
an unknown distribution $\D$ on $\X\times\Y$ with $\Y=\{-1,+1\}$, and we observe
a learning sample $S=\{(x_i,y_i)\}_{i=1}^m\sim\D^m$. We are given a set $\Hcal$ of
voters and a loss function $\ell:\Hcal\times(\X\times\Y)\to[0,1]$, which will be
the 0-1 loss $\ell(h,(x,y))=\ind{h(x)\neq y}$ unless stated otherwise. The true
risk and the empirical risk of a voter are
\[
  \RD(h)=\Esp{(x,y)\sim\D}{\ell(h,(x,y))}\qquad\text{and}\qquad
  \RS(h)=\frac1m\sum_{i=1}^m\ell(h,(x_i,y_i)).
\]
For voters with values in $\{-1,+1\}$ we will often use the identity
$\ind{h(x)\neq y}=\frac{1-y\,h(x)}{2}$, which turns the 0-1 loss into a linear
function of the prediction.

\begin{definition}{Gibbs classifier and majority vote}{gibbs}
Let $Q\in\M(\Hcal)$ be a distribution over $\Hcal$.
The \textbf{Gibbs classifier} $\GQ$ predicts the label of $x$ by drawing
$h\sim Q$ and returning $h(x)$. Its true and empirical risks are
\[
  \RD(\GQ)=\Esp{h\sim Q}{\RD(h)},\qquad \RS(\GQ)=\Esp{h\sim Q}{\RS(h)}.
\]
The \textbf{($Q$-weighted) majority vote} is the deterministic classifier
$\BQ(x)=\sign\big(\Esp{h\sim Q}{h(x)}\big)$, with the convention
$\sign(0)=-1$, so that a tie is counted as an error in the analysis. Its risk is
$\RD(\BQ)=\P_{(x,y)\sim\D}\big(\BQ(x)\neq y\big)$.
\end{definition}

The two predictors behave very differently. The Gibbs risk is \emph{linear} in
$Q$, which makes it easy to analyse, and it always lies between the risk of the
best and the risk of the worst voter. The risk of the majority vote is not
linear in $Q$, and it can be much smaller than the Gibbs risk: in the running
example of \cref{chap:intro}, the Gibbs risk of the forest is $19\%$ while its
vote only errs on $4.4\%$ of the test points. A convenient way to relate the two
is to exchange the expectations over the voters and over the data (Fubini's theorem):
\[
\RD(\GQ)=\Esp{(x,y)\sim\D}{W_Q(x,y)},\qquad\text{where}\qquad
W_Q(x,y)=\Esp{h\sim Q}{\ind{h(x)\neq y}}
\]
is the proportion, weighted by $Q$, of voters that err on $(x,y)$. The random
variable $W_Q$ will be the main character of \cref{chap:mv}. For now, the
important point is that the Gibbs risk is an average of the risks of single
voters, so our first task is to control the risk of one voter.

\subsection{Concentration inequalities}

All generalization bounds rely on the same basic fact: for a \emph{fixed}
voter $h$, chosen before seeing the data, the empirical risk $\RS(h)$ is an
average of $m$ independent random variables with values in $[0,1]$ and mean
$\RD(h)$, hence it concentrates around $\RD(h)$. The tools below quantify this
concentration. We start with the simplest of all.

\begin{lemma}{Markov's inequality}{markov}
Let $Z\geq 0$ be a random variable and $a>0$. Then $\P(Z\geq a)\leq \frac{\E Z}{a}$.
Equivalently, for any $\delta\in(0,1]$, with probability at least $1-\delta$, $Z\leq\frac{\E Z}{\delta}$.
\end{lemma}
\begin{proof}
Since $Z\geq0$, we have $a\ind{Z\geq a}\leq Z$; take expectations.
\end{proof}

Markov's inequality looks weak, but it becomes remarkably powerful when it is
applied to a well-chosen transformation of the variable of interest. Applied to
the exponential $e^{\lambda(p-\hat p)}$, it gives Hoeffding's inequality.

\begin{lemma}{Hoeffding's inequality \citep{hoeffding1963probability}}{hoeffding}
Let $Z_1,\dots,Z_m$ be i.i.d.\ random variables in $[0,1]$ with mean $p$, and
$\hat p=\frac1m\sum_i Z_i$. For any $\lambda\in\R$,
$\E\,e^{\lambda(p-\hat p)}\leq e^{\lambda^2/(8m)}$, and, for $\varepsilon>0$,
$\P(p-\hat p\geq\varepsilon)\leq e^{-2m\varepsilon^2}$.
\end{lemma}
\begin{proofsketch}
The first inequality is Hoeffding's lemma, applied to each of the $m$
independent terms $\frac{p-Z_i}{m}$, which have zero mean and range $\frac1m$.
For the second, Markov's inequality gives
$\P(p-\hat p\geq\varepsilon)\leq e^{-\lambda\varepsilon}\E e^{\lambda(p-\hat p)}\leq e^{-\lambda\varepsilon+\lambda^2/(8m)}$,
and the choice $\lambda=4m\varepsilon$ gives $e^{-2m\varepsilon^2}$.
\end{proofsketch}

Inverting Hoeffding's inequality, the true risk of a fixed voter satisfies
$\RD(h)\leq\RS(h)+\sqrt{\ln(1/\delta)/(2m)}$ with probability at least
$1-\delta$. In words, the empirical risk underestimates the true risk by more
than $\sqrt{\ln(1/\delta)/(2m)}$ with probability at most $\delta$. This is already
a certificate, but its width does not depend on the value of the risk, which is
wasteful when the risk is small: a Bernoulli variable of mean $0.01$ fluctuates
much less than a Bernoulli variable of mean $0.5$.

The sharpest statement for Bernoulli variables uses the \textbf{binary
Kullback--Leibler divergence}
\[
\kl(q\|p)=q\ln\frac{q}{p}+(1-q)\ln\frac{1-q}{1-p},\qquad q,p\in[0,1],
\]
which is the Kullback--Leibler divergence between two Bernoulli distributions of
parameters $q$ and $p$.

The binary kl has a very concrete meaning: it measures how \emph{surprising} it
is to observe a frequency $q$ of errors when the true error rate is $p$. Indeed,
the probability that $m$ independent coin flips of bias $p$ produce a frequency
close to $q$ behaves like $e^{-m\,\kl(q\|p)}$. This is Sanov's theorem for coins,
and the next lemma makes one side of it precise. A small kl means that the
observation is plausible under $p$, while a kl of $\varepsilon$ means that it has a
probability of order $e^{-m\varepsilon}$. Turning this around, the true risks $p$
compatible with an observed empirical risk $\hat q$ are those for which
$\kl(\hat q\|p)$ is small.

\begin{lemma}{Chernoff's bound and Maurer's lemma}{chernoff}
Let $Z_1,\dots,Z_m$ be i.i.d.\ random variables in $[0,1]$ with mean $p$ and
average $\hat p$. Chernoff's bound states that, for $\varepsilon\in(0,p)$,
$\P\big(\hat p\leq p-\varepsilon\big)\leq e^{-m\,\kl(p-\varepsilon\|p)}$.
Maurer's lemma \citep{maurer2004note} states that, for $m\geq 8$,
\begin{equation}\label{eq:maurer}
  \E\Big[e^{m\,\kl(\hat p\|p)}\Big]\leq \xi(m):=\sum_{k=0}^m\binom mk\Big(\frac km\Big)^k\Big(1-\frac km\Big)^{m-k},
\end{equation}
with $\sqrt m\leq\xi(m)\leq 2\sqrt m$; the first inequality in \eqref{eq:maurer} is an
equality when the $Z_i$ are Bernoulli variables.
\end{lemma}

The two statements say the same thing in two different forms. Chernoff's bound
controls a \emph{tail}: a large deviation of the empirical mean is exponentially
unlikely, at the rate given by the kl. Maurer's lemma controls an
\emph{exponential moment}: on average, $e^{m\kl(\hat p\|p)}$ is only of order
$\sqrt m$, although it could a priori be as large as $e^{m\ln 2}$. The first form
is the natural one for a single classifier; the second is the one that
PAC-Bayesian proofs need, because an expectation can be averaged over voters,
whereas a tail probability cannot.

\begin{proofsketch}
The key step in the second statement is that, for any convex function $f$ on
$[0,1]^m$, $\E f(Z_1,\dots,Z_m)\leq \E f(B_1,\dots,B_m)$ where the $B_i$ are
i.i.d.\ Bernoulli$(p)$ variables: a variable in $[0,1]$ is ``less spread'' than
a Bernoulli variable with the same mean. Since $q\mapsto e^{m\kl(q\|p)}$ is
convex, it suffices to treat the case $m\hat p\sim\mathrm{Bin}(m,p)$, for which
\[
\E e^{m\kl(\hat p\|p)}=\sum_k \binom mk p^k(1-p)^{m-k}\Big(\frac{k/m}{p}\Big)^{k}\Big(\frac{1-k/m}{1-p}\Big)^{m-k}=\xi(m),
\]
a quantity which, remarkably, does not depend on $p$; Stirling's formula gives
the bounds on $\xi(m)$.
\end{proofsketch}

\paragraph{Inverting the kl.} Both statements bound the kl between the empirical
and the true risk, whereas what we want is an explicit upper bound on the true
risk. We therefore need to invert the binary kl. For a fixed $q$, the function
$p\mapsto\kl(q\|p)$ is convex, vanishes at $p=q$ and increases on $[q,1]$ (see
\cref{fig:klinv}). The inequality $\kl(\hat q\|p)\leq\varepsilon$ therefore defines
an interval $[\kllow(\hat q,\varepsilon),\klup(\hat q,\varepsilon)]$ of admissible
values of $p$, with
\[
  \klup(\hat q,\varepsilon)=\sup\{p\in[\hat q,1]:\kl(\hat q\|p)\leq\varepsilon\},\qquad
  \kllow(\hat q,\varepsilon)=\inf\{p\in[0,\hat q]:\kl(\hat q\|p)\leq\varepsilon\}.
\]
These quantities have no closed form, but since the function is monotone on each
side of $\hat q$ they are computed to any precision by a few iterations of
bisection (\cref{algo:klinv}). In the rest of the notes, $\klup$ will be our
standard way of turning an empirical quantity into a certified upper bound.

\begin{figure}[t]
\centering
\includegraphics[width=0.62\linewidth]{kl_inverse}
\caption{The function $p\mapsto\kl(\hat q\|p)$ for $\hat q=0.1$. For
$\varepsilon=0.05$, the values of $p$ compatible with $\kl(\hat q\|p)\leq\varepsilon$ form
an interval whose upper end $\klup(\hat q,\varepsilon)=0.220$ is smaller than the
Pinsker relaxation $\hat q+\sqrt{\varepsilon/2}=0.258$.}
\label{fig:klinv}
\end{figure}

\begin{algorithm}[t]
\caption{Upper inversion of the binary kl by bisection}\label{algo:klinv}
\KwIn{empirical value $\hat q\in[0,1]$, budget $\varepsilon\geq 0$, tolerance $\tau$}
$a\leftarrow\hat q$, $b\leftarrow 1$\;
\While{$b-a>\tau$}{
  $c\leftarrow(a+b)/2$\;
  \lIf{$\kl(\hat q\|c)>\varepsilon$}{$b\leftarrow c$}\lElse{$a\leftarrow c$}
}
\KwOut{$b$ \tcp*{$\geq\klup(\hat q,\varepsilon)$, within $\tau$}}
\end{algorithm}

\begin{example}[Inverting the kl by hand]
Let $\hat q=0.1$ and $\varepsilon=0.05$. We know that $\klup(0.1,0.05)\geq0.1$.
Since $\kl(0.1\|0.2)\approx0.037<0.05$ and $\kl(0.1\|0.25)\approx0.072>0.05$, the
answer lies in $[0.2,0.25]$; the midpoint gives $\kl(0.1\|0.225)\approx0.053>0.05$,
so the answer is in $[0.2,0.225]$, and so on. After twenty steps the interval
has length $10^{-6}$ and we find $\klup(0.1,0.05)\approx0.220$.
\end{example}

Bisection gives numbers, but not formulas. To understand how kl-based bounds
behave, and to derive bounds with a closed form, we also need explicit (if
looser) upper bounds on $\klup$; the next two inequalities provide them.

\begin{lemma}{Pinsker-type inequalities}{pinsker}
For $q,p\in[0,1]$, $\kl(q\|p)\geq 2(q-p)^2$ (Pinsker's inequality), hence
$\klup(q,\varepsilon)\leq q+\sqrt{\varepsilon/2}$. Moreover, if $q<p$, then
$\kl(q\|p)\geq\frac{(p-q)^2}{2p}$ (refined Pinsker inequality), hence
$\klup(q,\varepsilon)\leq q+\sqrt{2q\varepsilon}+2\varepsilon$.
\end{lemma}

Pinsker's inequality recovers a Hoeffding-type width $\sqrt{\varepsilon/2}$. The
refined inequality explains why it is so much better for small risks: when $q$ is close to zero, the width of the
confidence interval is of order $\varepsilon$ instead of $\sqrt\varepsilon$, i.e.\ of order
$1/m$ instead of $1/\sqrt m$ when $\varepsilon\propto1/m$. This ``fast rate'' is
visible on the left panel of \cref{fig:concentration}: for an empirical risk of
$1\%$ and $m=1000$ examples, the kl interval is four times narrower than the
Hoeffding interval. The right panel checks by simulation that both intervals
are valid: their failure frequency always stays below $\delta$, and usually well
below, as these inequalities are worst-case statements.

\begin{running}
Suppose that we had fixed the forest before seeing the $540$ test digits, and
that we use them only to certify it: this is the classical \emph{test-set bound}.
The vote errs on $4.4\%$ of them; with $\delta=0.05$, Hoeffding's inequality
certifies a true risk below $0.044+\sqrt{\ln 20/1080}=0.097$, and the kl
inversion below $0.069$. For an average tree ($19\%$ of test errors), the two
certificates are $0.243$ and $0.234$: the kl only makes a difference for small
risks. These numbers will serve as a reference. Our goal is to obtain comparable
guarantees \emph{without} setting aside a test sample, from the training data
alone.
\end{running}

\begin{figure}[t]
\centering
\includegraphics[width=\linewidth]{concentration}
\caption{Left: width $p-\hat p$ of the upper confidence bound on the risk of a
fixed voter given by Hoeffding's inequality (constant) and by the kl inversion,
as a function of the empirical risk, for $m=100$ and $m=1000$. Right: frequency,
over $4000$ simulated samples of size $100$, of the event ``the true risk is
above the bound''; both bounds are valid.}
\label{fig:concentration}
\end{figure}

\paragraph{Second moments.} The inequalities above only use the fact that the
variables are bounded. Majority votes will require a finer tool, which also
uses the \emph{variance} of a random variable, and not only its mean.

\begin{lemma}{Chebyshev--Cantelli inequality}{cantelli}
Let $Z$ be a random variable with mean $\mu$ and variance $v$. For every $a>0$,
$\P(Z-\mu\geq a)\leq\frac{v}{v+a^2}$ and $\P(Z-\mu\leq-a)\leq\frac{v}{v+a^2}$.
\end{lemma}
\begin{proof}
For any $t\geq0$, $\{Z-\mu\geq a\}\subseteq\{(Z-\mu+t)^2\geq(a+t)^2\}$, so by Markov's
inequality $\P(Z-\mu\geq a)\leq\frac{\E(Z-\mu+t)^2}{(a+t)^2}=\frac{v+t^2}{(a+t)^2}$. The
right-hand side is minimal for $t=v/a$, where it equals $\frac{v}{v+a^2}$. Applying
the result to $-Z$ gives the lower tail.
\end{proof}

In words, a variable whose variance is small compared with $a^2$ rarely
deviates from its mean by more than $a$, whatever its distribution. Note also
the structure of the proof, which we will meet again in \cref{chap:mv}: we control a tail probability by the expectation of a well-chosen non-negative
function, and then optimize over a free parameter of this function.

\subsection{The Kullback--Leibler divergence and the change of measure}

Everything so far concerns a voter fixed \emph{before} seeing the data. Learning
is precisely the opposite: the algorithm looks at the sample and chooses. To
account for this choice, PAC-Bayes compares the distribution $Q$ that has been
chosen with a distribution $P$ that was fixed in advance. We need two things: a
way to measure how far $Q$ is from $P$, and a way to transfer statements that are
true under $P$ into statements under $Q$. The Kullback--Leibler divergence and
the change of measure inequality provide exactly these two things.

\begin{definition}{Kullback--Leibler divergence}{kl}
For $Q,P\in\M(\Hcal)$ such that $Q$ is absolutely continuous w.r.t.\ $P$,
\[
  \KL(Q\|P)=\Esp{h\sim Q}{\ln\frac{dQ}{dP}(h)}
  \quad\Big(=\sum_{h\in\Hcal}Q(h)\ln\frac{Q(h)}{P(h)}\ \text{ if $\Hcal$ is finite}\Big),
\]
and $\KL(Q\|P)=+\infty$ otherwise.
\end{definition}

The KL divergence is always non-negative and it vanishes only when $Q=P$; this
is Gibbs' inequality, a consequence of Jensen's inequality applied to the
concave logarithm. It is not symmetric, and it is not a distance, but it
behaves as a measure of how far the posterior has moved away from the prior.
A few examples help to get a feeling for its values.

On a finite set of $n$ voters with the uniform prior, it takes the simple form
$\KL(Q\|P)=\ln n-H(Q)$, where $H(Q)=-\sum_hQ(h)\ln Q(h)$ is the Shannon entropy
of $Q$. The KL is then zero for the uniform posterior and maximal, equal to
$\ln n$, for a Dirac on a single voter; if $Q$ is uniform on $k$ of the $n$
voters, $\KL(Q\|P)=\ln\frac nk$: concentrating on a tenth of the voters costs
$\ln 10$ nats. For product distributions the KL divergence
adds up, $\KL(Q_1\otimes Q_2\|P_1\otimes P_2)=\KL(Q_1\|P_1)+\KL(Q_2\|P_2)$; in
particular $\KL(Q^2\|P^2)=2\KL(Q\|P)$ where $Q^2=Q\otimes Q$ is the distribution
of a pair of independent voters, an identity which will be used for majority
votes. Finally, for Gaussian distributions with the same covariance,
$\KL\big(\mathcal N(\mu_1,\sigma^2 I_d)\,\|\,\mathcal N(\mu_0,\sigma^2 I_d)\big)=\frac{\|\mu_1-\mu_0\|^2}{2\sigma^2}$.
These properties are proved in the tutorial sheet.

We now have a way to measure how far $Q$ is from $P$; it remains to transfer
statements from one to the other. The following inequality, the heart of every
PAC-Bayesian proof, does exactly that. It allows
us to replace an expectation under the data-dependent posterior $Q$ by an
expectation under the data-independent prior $P$, at the price of $\KL(Q\|P)$.

\begin{lemma}{Change of measure \citep{donsker1975asymptotic,csiszar1975divergence}}{dv}
For any measurable function $\phi:\Hcal\to\R$ and any $P\in\M(\Hcal)$,
\begin{equation}\label{eq:dv}
  \forall Q\in\M(\Hcal),\qquad
  \Esp{h\sim Q}{\phi(h)}\ \leq\ \KL(Q\|P)+\ln\Esp{h\sim P}{e^{\phi(h)}}.
\end{equation}
Moreover, when $\E_P e^{\phi}<\infty$, equality holds for the
\emph{Gibbs distribution} $Q_\phi(h)=\frac{e^{\phi(h)}P(h)}{\E_{P}e^{\phi}}$, so that
\begin{equation}\label{eq:gibbs_variational}
  \ln\Esp{h\sim P}{e^{\phi(h)}}=\sup_{Q\in\M(\Hcal)}\Big\{\Esp{h\sim Q}{\phi(h)}-\KL(Q\|P)\Big\}.
\end{equation}
\end{lemma}

\begin{proof}
Assume $\KL(Q\|P)<\infty$ (otherwise there is nothing to prove) and let
$Q_\phi$ be the Gibbs distribution. By definition of $Q_\phi$,
\[
  \KL(Q\|Q_\phi)=\Esp{h\sim Q}{\ln\frac{Q(h)}{P(h)}-\phi(h)+\ln\E_{P}e^{\phi}}
  =\KL(Q\|P)-\E_Q\phi+\ln\E_P e^{\phi}.
\]
Since $\KL(Q\|Q_\phi)\geq0$, we obtain \eqref{eq:dv}, with equality if and only if $Q=Q_\phi$.
\end{proof}

In plain words, the average of $\phi$ under \emph{any} posterior is controlled
by an exponential average under the prior, which does not depend on $Q$, plus a
penalty that grows as $Q$ moves away from $P$. Since the inequality holds
for all $Q$ at once, the posterior may be chosen after the fact. The proof shows
even more than the inequality: the gap in \eqref{eq:dv} is exactly
$\KL(Q\|Q_\phi)$, the divergence between $Q$ and the Gibbs distribution.
\Cref{fig:change_of_measure} illustrates this on a small example with $25$
voters: along the segment joining an arbitrary posterior $Q$ to $Q_\phi$, the
quantity $\E_Q\phi-\KL(Q\|P)$ increases until it reaches the value
$\ln\E_Pe^\phi$ at $Q_\phi$.

\begin{figure}[t]
\centering
\includegraphics[width=\linewidth]{change_of_measure}
\caption{The change of measure on $25$ voters. Left: a prior $P$, a function $\phi$
(black dots), an arbitrary posterior $Q$ and the Gibbs distribution $Q_\phi\propto Pe^\phi$,
which puts its mass on the voters where $\phi$ is large while remaining close to $P$.
Right: value of $\E_{Q_t}\phi-\KL(Q_t\|P)$ along $Q_t=(1-t)Q+tQ_\phi$; it never exceeds
$\ln\E_Pe^\phi$ and reaches it at $Q_\phi$.}
\label{fig:change_of_measure}
\end{figure}

\begin{intuition}
Equation~\eqref{eq:gibbs_variational} says that $\ln\E_P e^\phi$ is the
maximal value of a ``fit minus complexity'' criterion $\E_Q\phi-\KL(Q\|P)$, and
that the best posterior is obtained by \emph{exponentially reweighting} the
prior. This variational principle is the origin of the \emph{Gibbs posterior}
(\cref{sec:gibbs_posterior}) and of the link between PAC-Bayes and Bayesian inference.
\end{intuition}

\subsection{Warm-up: the Occam bound}\label{sec:occam}

Before combining the two tools in full generality, let us see how far the
concentration inequalities alone can take us when the learner selects one
classifier among many. Consider a \emph{countable} class $\Hcal$ and a
prior $P$ on $\Hcal$, which can be thought of as a ``description length'':
$\ln\frac1{P(h)}$ is the number of nats needed to describe $h$ with a code
adapted to $P$.

\begin{theorem}{Occam's razor bound \citep{blumer1987occam,langford2005tutorial}}{occam}
For any countable $\Hcal$, any prior $P$ with $P(h)>0$ for all $h$, and any
$\delta\in(0,1]$, with probability at least $1-\delta$ over $S\sim\D^m$,
\begin{gather*}
  \forall h\in\Hcal,\qquad \RD(h)\leq\klup\Big(\RS(h),\frac{\ln\frac{1}{P(h)}+\ln\frac1\delta}{m}\Big),\\
  \text{hence}\qquad \RD(h)\leq\RS(h)+\sqrt{\frac{\ln\frac1{P(h)}+\ln\frac1\delta}{2m}}.
\end{gather*}
\end{theorem}

\begin{proof}
For a fixed $h$, $m\RS(h)$ is a sum of $m$ i.i.d.\ Bernoulli variables of mean
$\RD(h)$. By Chernoff's bound (lower tail of $\RS(h)$), the event
\[
\big\{\RS(h)<\RD(h)\text{ and }\kl(\RS(h)\|\RD(h))>\varepsilon_h\big\}=\big\{\RD(h)>\klup(\RS(h),\varepsilon_h)\big\}
\]
has probability at most $e^{-m\varepsilon_h}=\delta P(h)$ when
$\varepsilon_h=\frac1m\ln\frac{1}{\delta P(h)}$. A union bound over $h\in\Hcal$ gives
a total failure probability at most $\sum_h\delta P(h)=\delta$. The second
statement follows from Pinsker's inequality.
\end{proof}

The theorem certifies every hypothesis of the class at once, so it remains valid
for the one the learner finally picks. Its proof shows how the prior works: it \emph{distributes the confidence budget}
$\delta$ over the hypotheses, each $h$ receiving the share $\delta P(h)$. A
hypothesis that was \emph{a priori} likely is cheap, and an unlikely one is
expensive. With the uniform prior on a finite class, $\ln\frac1{P(h)}=\ln|\Hcal|$
and we recover the classical bound for finite classes seen in the first part of
the course. As shown in \cref{fig:occam}, the price of a large finite class is
surprisingly small: with $m=1000$ examples and an empirical risk of $5\%$,
going from $10$ to $10^{12}$ candidate classifiers only moves the bound from
about $0.08$ to $0.12$.

\begin{figure}[t]
\centering
\includegraphics[width=0.66\linewidth]{occam}
\caption{The Occam bound (kl form, $\delta=0.05$, $m=1000$) as a function of the
size of the class, for three values of the empirical risk. The complexity only
grows like $\ln|\Hcal|$.}
\label{fig:occam}
\end{figure}

The Occam bound already has the PAC-Bayesian flavour, but it has two limits. It
only certifies a \emph{single} hypothesis --- in the language of the next
sections, a Dirac posterior, for which $\KL(\delta_h\|P)=\ln\frac1{P(h)}$ --- and
it is useless for uncountable classes, where $P(h)=0$ for every $h$. PAC-Bayes
removes both limits by considering \emph{averages} of hypotheses: averaging
over a posterior which puts a large mass around a good hypothesis costs much
less than pinpointing this hypothesis exactly. Making this idea precise, with
the change of measure in place of the union bound, is the purpose of the next
section.

\begin{keypoints}
For a fixed voter, the empirical risk concentrates around the true risk; the
binary kl gives the sharpest statement, and its inversion $\klup$, computed by
bisection, turns an empirical risk into a certified upper bound, with a fast
rate when the risk is small. The change of measure inequality
$\E_Q\phi\leq\KL(Q\|P)+\ln\E_Pe^\phi$ transfers a statement about a
data-independent prior to any posterior, at the price $\KL(Q\|P)$, and it is an
equality for the Gibbs distribution. The Occam bound is the special case of
Dirac posteriors on a countable class; its price is $\ln\frac1{P(h)}$.
\end{keypoints}

\begin{quickchecks}
\item Why is the kl confidence interval narrower than Hoeffding's when the empirical risk is small?
\item What is $\KL(Q\|P)$ when $P$ is uniform on $n$ voters and $Q$ is uniform on $k$ of them?
\item For which posterior is the change of measure inequality an equality?
\item In the Occam bound, what plays the role of the complexity of the chosen hypothesis?
\end{quickchecks}

\begin{exercises}
\begin{exercise}
Compute $\kl(0.1\|0.2)$ and $\kl(0.2\|0.1)$. Is the binary kl symmetric?
\end{exercise}
\begin{exercise}
Show that the Gibbs distribution $Q_\phi$ of Lemma~\ref{lem:dv} is the only posterior
for which the change of measure is an equality.
\end{exercise}
\begin{exercise}
Using the Occam bound in Pinsker form, how many examples are needed to certify a
risk below $0.15$ for a classifier with empirical risk $0.1$ chosen among $10^6$
candidates, with $\delta=0.05$?
\end{exercise}
\begin{exercise}
Using Pinsker's inequality, show that the kl interval of \cref{fig:concentration}
is never wider than the Hoeffding interval. Using notebook~1, determine for which
empirical risks it is at least twice narrower when $m=1000$.
\end{exercise}
\begin{exercise}
\textbf{(Comprehension)} True or false? Justify briefly. (a)~If $P$ and $Q$ are
two distributions on a finite set of voters, $\KL(Q\|P)$ is finite. (b)~If $P$ is
uniform on $n$ voters, then $\KL(Q\|P)\leq\ln n$ for every posterior $Q$.
(c)~In the Occam bound with the uniform prior, doubling the number of candidate
classifiers doubles the complexity term $\ln\frac1{P(h)}$.
\end{exercise}
\begin{exercise}
\textbf{(Application)} Let $P$ be uniform on $n=4$ voters. Compute $\KL(Q\|P)$
for $Q=(\frac12,\frac14,\frac14,0)$, for $Q'=(0.7,0.1,0.1,0.1)$ and for the Dirac
$\delta_{h_1}$, both from the definition and with the formula $\ln n-H(Q)$.
Compute also $\KL(P\|Q')$ and $\KL(P\|Q)$, and comment.
\end{exercise}
\begin{exercise}
\textbf{(Application)} Check the two test-set certificates of the running
example of this section: a vote with $4.4\%$ of errors on $540$ test points,
certified by Hoeffding's inequality ($0.097$) and by the kl inversion ($0.069$) with
$\delta=0.05$. Recompute both certificates with $\delta=0.01$, and comment on the
cost of the extra confidence.
\end{exercise}
\begin{exercise}
\textbf{(Proof)} Prove Markov's inequality (Lemma~\ref{lem:markov}) in detail, then
deduce Chebyshev's inequality: for a random variable $Z$ with mean $\mu$ and
variance $v$, $\P(|Z-\mu|\geq a)\leq\frac{v}{a^2}$ for every $a>0$. Compare it with the
one-sided Cantelli inequality (Lemma~\ref{lem:cantelli}) when $a^2=v$.
\end{exercise}
\begin{exercise}
\textbf{(Proof)} Admitting Hoeffding's lemma
$\E\,e^{\lambda(p-\hat p)}\leq e^{\lambda^2/(8m)}$ (Lemma~\ref{lem:hoeffding}), prove the tail bound
$\P(p-\hat p\geq\varepsilon)\leq e^{-2m\varepsilon^2}$ by optimizing over $\lambda>0$, and deduce that
$p\leq\hat p+\sqrt{\ln(1/\delta)/(2m)}$ with probability at least $1-\delta$.
\end{exercise}
\begin{exercise}
\textbf{(Proof)} Let $Q_1,P_1$ be distributions on a finite set $\Hcal_1$ and
$Q_2,P_2$ distributions on a finite set $\Hcal_2$, with $P_1,P_2$ positive. Show that
$\KL(Q_1\otimes Q_2\|P_1\otimes P_2)=\KL(Q_1\|P_1)+\KL(Q_2\|P_2)$, and deduce that
$\KL(Q^2\|P^2)=2\KL(Q\|P)$, the identity used for the tandem bound. Show also that
$\KL(\delta_h\|P)=\ln\frac1{P(h)}$ for a Dirac posterior.
\end{exercise}
\end{exercises}
